\documentclass[10pt,a4paper]{article}
\usepackage[margin=20mm,headheight=15pt,headsep=7mm,footskip=10mm]{geometry}
\usepackage{fontspec,amsmath,amssymb,mathtools,microtype}
\setmainfont{lmroman10-regular.otf}[BoldFont=lmroman10-bold.otf,ItalicFont=lmroman10-italic.otf,BoldItalicFont=lmroman10-bolditalic.otf]
\setsansfont{lmsans10-regular.otf}[BoldFont=lmsans10-bold.otf,ItalicFont=lmsans10-oblique.otf,BoldItalicFont=lmsans10-boldoblique.otf]
\setmonofont[Scale=.86]{lmmono10-regular.otf}[ItalicFont=lmmono10-italic.otf]
\usepackage{ragged2e}\usepackage{xcolor,booktabs,longtable,array,graphicx,enumitem,fancyhdr,needspace,titlesec,hyperref,xurl}
\definecolor{AtlasNavy}{HTML}{18334C}\definecolor{AtlasTeal}{HTML}{116D73}\definecolor{AtlasGray}{HTML}{52616D}
\hypersetup{colorlinks=true,linkcolor=AtlasTeal,urlcolor=AtlasTeal,pdfcreator={Complexity Atlas},unicode=true}
\urlstyle{same}\setlength{\parindent}{0pt}\setlength{\parskip}{4.8pt}\setlength{\emergencystretch}{3em}
\renewcommand{\baselinestretch}{1.08}\setcounter{secnumdepth}{0}\setlist{nosep,leftmargin=1.5em}
\titleformat{\section}{\Large\bfseries\sffamily\color{AtlasNavy}\raggedright}{}{0em}{}
\titleformat{\subsection}{\large\bfseries\sffamily\color{AtlasTeal}\raggedright}{}{0em}{}
\titleformat{\subsubsection}{\normalsize\bfseries\color{AtlasNavy}\raggedright}{}{0em}{}
\titlespacing*{\section}{0pt}{14pt}{6pt}\titlespacing*{\subsection}{0pt}{10pt}{5pt}
\pagestyle{fancy}\fancyhf{}\fancyhead[L]{\small\sffamily\color{AtlasGray}COMPLEXITY ATLAS}
\fancyhead[R]{\small\sffamily\color{AtlasGray}ENGLISH EDITION}\fancyfoot[C]{\small\thepage}\renewcommand{\headrulewidth}{.3pt}
\newcommand{\AtlasMath}[1]{\par\begingroup\centering\setbox0=\hbox{$\displaystyle #1$}\ifdim\wd0>1.4\linewidth\typeout{ATLAS WIDE FORMULA: \the\wd0}\fi\ifdim\wd0>\linewidth\resizebox{\linewidth}{!}{\box0}\else\box0\fi\par\endgroup}
\newcommand{\AtlasInlineMath}[1]{\begingroup\setbox0=\hbox{$#1$}\ifdim\wd0>.4\linewidth\endgroup\AtlasMath{#1}\else\endgroup$#1$\fi}
\pdfstringdefDisableCommands{\def\AtlasInlineMath#1{#1}}
\DeclareMathOperator{\polylog}{polylog}\DeclareMathOperator{\dist}{dist}\DeclareMathOperator{\diam}{diam}
\DeclareMathOperator{\rank}{rank}\DeclareMathOperator{\tr}{tr}\DeclareMathOperator{\diag}{diag}
\DeclareMathOperator{\prox}{prox}\DeclareMathOperator{\KL}{KL}\DeclareMathOperator{\Var}{Var}
\providecommand{\eps}{\varepsilon}\providecommand{\R}{\mathbb R}\providecommand{\E}{\mathbb E}
\providecommand{\mathbbm}[1]{\mathbf{#1}}\providecommand{\mathds}[1]{\mathbf{#1}}
\providecommand{\llbracket}{[\![}\providecommand{\rrbracket}{]\!]}
\begin{document}
\raggedright

{\LARGE\bfseries\sffamily\color{AtlasNavy} Perturbed block multiplication: the fixed Noisy Power Method\par}\vspace{5pt}
\textbf{Setting ID.} \texttt{MX\_\allowbreak{}S\_\allowbreak{}noisy\_\allowbreak{}power}

Matrix computation

The fixed Noisy Power Method returns a subspace of width \AtlasInlineMath{b_{\mathrm{blk}}} that approximately contains the true leading \AtlasInlineMath{k}-dimensional singular subspace.

\textbf{Status.} Bounds missing

\textbf{TASK tags.} task family: Matrix computation

\textbf{MODEL tags.} spectral structure: Explicit matrix class; spectral structure: Symmetric singular gap

\textbf{ORACLE tags.} feedback quality: Bounded adversarial; interface family: Causal bounded block noise; operation: Noisy block matvec; public information: Hidden objects explicitly declared; public information: Public inputs explicitly declared

\textbf{ALGORITHMS tags.} adaptivity: Causal protocol

\textbf{OUTPUT tags.} certificate: No additional output certificate required; object: Orthonormal block

\textbf{METRIC tags.} criterion: Subspace containment error

\textbf{GUARANTEE tags.} randomness: Algorithm private seed; time uniformity: Fixed time or accuracy

\textbf{RESOURCE tags.} accounting: Joint coordinates retained; measure: N block; measure: N vec

\needspace{7\baselineskip}\section{1. Setting specification}
\needspace{4\baselineskip}\subsection{Task}
The fixed Noisy Power Method returns a subspace of width \AtlasInlineMath{b_{\mathrm{blk}}} that approximately contains the true leading \AtlasInlineMath{k}-dimensional singular subspace.

\needspace{4\baselineskip}\subsection{Admissible model}
Assume \AtlasInlineMath{A=A^{\mathsf T}\in\mathbb{R}^{d\times d}}, \AtlasInlineMath{\sigma_1(A)\le\Lambda}, and \AtlasInlineMath{\sigma_k(A)-\sigma_{k+1}(A)\ge g>0}. Let \AtlasInlineMath{U\in\mathbb{R}^{d\times k}} be an orthonormal basis for the leading singular subspace, with \AtlasInlineMath{1\le k\le b_{\mathrm{blk}}\le d}; append zero singular values at the end when needed.

\needspace{4\baselineskip}\subsection{Oracle interface}
Inputs \AtlasInlineMath{d,k,b_{\mathrm{blk}},\Lambda,g,\nu_{\mathrm{tot}},\nu_{\parallel},\varepsilon,\delta} are public and \AtlasInlineMath{A,U} is hidden. For a block \AtlasInlineMath{X\in\mathbb{R}^{d\times b_{\mathrm{blk}}}} with orthonormal columns, the response is \AtlasInlineMath{AX+G}. The oracle may choose \AtlasInlineMath{G} from the history, but \AtlasInlineMath{\|G\|_{\mathrm{op}}\le\nu_{\mathrm{tot}}} and \AtlasInlineMath{\|U^{\mathsf T}G\|_{\mathrm{op}}\le\nu_{\parallel}} must hold on every realization at every step.

\needspace{4\baselineskip}\subsection{Output}
Return the fixed algorithm’s final orthogonal block \AtlasInlineMath{X_L} with \AtlasInlineMath{b_{\mathrm{blk}}} columns. Compression to rank k is not required.

\needspace{4\baselineskip}\subsection{Loss or accuracy criterion}
The base error is \AtlasInlineMath{\|(I-X_LX_L^{\mathsf T})U\|_{\mathrm{op}}}, with success threshold \AtlasInlineMath{\varepsilon}. It measures how well the output subspace contains the true target subspace.

\needspace{4\baselineskip}\subsection{Resource accounting}
Retain joint resources \AtlasInlineMath{(N_{\mathrm{block}},N_{\mathrm{vec}})}. A block call of width \AtlasInlineMath{b_{\mathrm{blk}}} increments the first coordinate by one and the second by \AtlasInlineMath{b_{\mathrm{blk}}}, giving \AtlasInlineMath{N_{\mathrm{block}}=L} and \AtlasInlineMath{N_{\mathrm{vec}}=b_{\mathrm{blk}}L}.

The second term converts block products to per-column cost. The actual interface remains an orthogonal block product with a joint perturbation constraint. Public parameters determine the iteration count in advance; orthogonalization, storage, and internal arithmetic are free.

\needspace{4\baselineskip}\subsection{Algorithm views and success guarantees}
Use the fixed Noisy Power Method with random orthonormal initialization and update \AtlasInlineMath{X_{\ell}=\operatorname{orth}(AX_{\ell-1}+G_{\ell})}. Orthogonalization is exact internal computation. The upper bound evaluates this algorithm, not an infimum over all algorithms.

For every admissible matrix and every causal perturbation rule satisfying the per-step norm bounds, probability is over random initialization alone. The error criterion holds with probability at least \AtlasInlineMath{1-\delta}. Public parameters determine the iteration count in advance.

\textbf{Budget convention.} fixed worst-case joint budget; no expectation substituted

\needspace{7\baselineskip}\section{2. Notation and parameters}
\AtlasInlineMath{d}: The order of the square input matrix and the dimension of each query vector. \textbf{Domain:} \AtlasInlineMath{d\in\mathbb N,\ d\ge1}.

\AtlasInlineMath{A\in\mathbb R^{d\times d}}

\AtlasInlineMath{k}: Dimension of the leading singular subspace to be contained \textbf{Domain:} \AtlasInlineMath{k\in\mathbb N,\ 1\le k\le b_{\mathrm{blk}}}.

\AtlasInlineMath{U\in\mathbb R^{d\times k},\ U^{\mathsf T}U=I_k}

\AtlasInlineMath{b_{\mathrm{blk}}}: The orthogonal-block width in every query and the final output. \textbf{Domain:} \AtlasInlineMath{b_{\mathrm{blk}}\in\mathbb N,\ k\le b_{\mathrm{blk}}\le d}.

\AtlasInlineMath{X_\ell\in\mathbb R^{d\times b_{\mathrm{blk}}}}

\AtlasInlineMath{\Lambda}: The public spectral-norm bound. \textbf{Domain:} \AtlasInlineMath{\Lambda>0}.

\AtlasInlineMath{\sigma_1(A)\le\Lambda}

\AtlasInlineMath{g}: A lower bound on the absolute singular-value gap at the target boundary. \textbf{Domain:} \AtlasInlineMath{0<g\le\Lambda}.

\AtlasInlineMath{\sigma_k(A)-\sigma_{k+1}(A)\ge g}

\AtlasInlineMath{\nu_{\mathrm{tot}}}: Spectral-norm bound on the complete perturbation block at each step \textbf{Domain:} \AtlasInlineMath{\nu_{\mathrm{tot}}\ge0}.

\AtlasInlineMath{\|G_\ell\|_{\rm op}\le\nu_{\mathrm{tot}}}

\AtlasInlineMath{\nu_\parallel}: A spectral-norm bound on each perturbation’s projection onto the true target subspace. \textbf{Domain:} \AtlasInlineMath{\nu_\parallel\ge0}.

\AtlasInlineMath{\|U^{\mathsf T}G_\ell\|_{\rm op}\le\nu_\parallel}

\AtlasInlineMath{\varepsilon}: Operator-norm tolerance for the portion of the target subspace outside the output space. \textbf{Domain:} \AtlasInlineMath{\varepsilon>0}.

\AtlasInlineMath{\|(I_d-X_LX_L^{\mathsf T})U\|_{\rm op}\le\varepsilon}

\AtlasInlineMath{\delta}: Allowed failure probability for each fixed admissible input \textbf{Domain:} \AtlasInlineMath{0<\delta<1/2}.

Success probability is at least \AtlasInlineMath{1-\delta}.

\AtlasInlineMath{A}: The hidden real symmetric input matrix. \textbf{Domain:} \AtlasInlineMath{\mathbb R^{d\times d}}.

\AtlasInlineMath{\sigma_i(A)}: Input singular values in nonincreasing order, padded with zeros at the end. \textbf{Domain:} \AtlasInlineMath{1\le i\le d+1}.

\AtlasInlineMath{\sigma_1(A)\ge\cdots\ge\sigma_d(A)\ge\sigma_{d+1}(A)=0}

\AtlasInlineMath{U}: An orthonormal basis for the true leading k-dimensional singular subspace. \textbf{Domain:} \AtlasInlineMath{\mathbb R^{d\times k}}.

\AtlasInlineMath{U^{\mathsf T}U=I_k}

\AtlasInlineMath{X_\ell}: Orthonormal column block retained at step \AtlasInlineMath{\ell} \textbf{Domain:} \AtlasInlineMath{\mathbb R^{d\times b_{\rm blk}}}.

\AtlasInlineMath{X_\ell^{\mathsf T}X_\ell=I_{b_{\rm blk}}}

\AtlasInlineMath{G_\ell}: Additive error block at step \AtlasInlineMath{\ell}, generated by the causal disturbance rule \textbf{Domain:} \AtlasInlineMath{\mathbb R^{d\times b_{\rm blk}}}.

\AtlasInlineMath{\operatorname{orth}}: Specified exact orthogonalization of a given column space. \textbf{Domain:} Used by the fixed Noisy Power Method.

\AtlasInlineMath{L}: Number of iterations determined by public parameters \textbf{Domain:} \AtlasInlineMath{\mathbb N_0}.

\AtlasInlineMath{\tau}: A tunable parameter in the random orthogonal-initialization tail bound. \textbf{Domain:} \AtlasInlineMath{\tau>1}.

\AtlasInlineMath{c_0,c_1}: Two universal positive constants in the random-initialization tail bound \textbf{Domain:} A fixed constant.

\AtlasInlineMath{\tau^{-c_0(b_{\rm blk}+1-k)}+e^{-c_1d}}

\AtlasInlineMath{N_{\mathrm{block}}}: The total native block-interface calls in the algorithm. \textbf{Domain:} \AtlasInlineMath{\mathbb N_0}.

\AtlasInlineMath{N_{\rm block}=L}

\AtlasInlineMath{N_{\mathrm{vec}}}: Cost coordinate that charges each block query by its number of columns \textbf{Domain:} This cost coordinate does not grant a separate single-column query interface.

\AtlasInlineMath{N_{\rm vec}=b_{\rm blk}N_{\rm block}}

\AtlasInlineMath{I_j}: Identity matrix of order \AtlasInlineMath{j} \textbf{Domain:} \AtlasInlineMath{\mathbb R^{j\times j}}.

\AtlasInlineMath{\|z\|_2}: Euclidean norm of a vector \textbf{Domain:} Finite-dimensional real or complex vector.

\AtlasInlineMath{\|z\|_2=(\sum_i|z_i|^2)^{1/2}}

\AtlasInlineMath{\|M\|_{\mathrm{op}}}: The matrix operator norm induced by the Euclidean norm. \textbf{Domain:} Any compatible matrix.

\AtlasInlineMath{\|M\|_{\mathrm{op}}=\max_{\|z\|_2=1}\|Mz\|_2=\sigma_1(M)}

\AtlasInlineMath{c,C}: Positive universal constants in the corresponding bound. Identically named constants in different theorems need not have the same value. \textbf{Domain:} Independent of the dimension, accuracy, and probability parameters tracked by this bound.

\AtlasInlineMath{\log}: Natural logarithm \textbf{Domain:} Positive real argument.

\AtlasInlineMath{\log x=\ln x}

\AtlasInlineMath{O,\Omega,\Theta}: Upper, lower, and two-sided bounds up to constants in the same applicable regime \textbf{Domain:} Hidden constants have only the parameter dependence explicitly permitted by the associated bound.

\needspace{12\baselineskip}\section{3. Complexity bracket}
\needspace{4\baselineskip}\subsection{Compared quantities}
\needspace{4\baselineskip}\subsubsection{Q1 - source scoped metric resource statement}
\textbf{Metric contract.}

\AtlasMath{\Pr(\|(I_d-X_LX_L^{\mathsf T})U\|_{\rm op}\le\varepsilon)\ge1-\delta}
\textbf{Resource contract.} vector: N\_block; N\_vec; projection: N\_vec=b\_blk*N\_block. N\_vec charges each column in a block call, while information access is still defined by the orthogonal-block oracle.

Subspace-containment error and joint-call guarantee for the fixed noisy power method. There is no minimization over all algorithms and no additional weighting of the resource vector.

\needspace{4\baselineskip}\subsection{R1 - The perturbation threshold and initialization-confidence window.}
\textbf{Region.} Choose \AtlasInlineMath{\tau>1} so that \AtlasInlineMath{\tau^{-c_0(b_{\mathrm{blk}}+1-k)}+e^{-c_1d}\le\delta}. Require \AtlasInlineMath{0<\varepsilon<1/2}, \AtlasInlineMath{5\nu_{\mathrm{tot}}\le\varepsilon g}, and \AtlasInlineMath{5\nu_{\parallel}\le g(\sqrt{b_{\mathrm{blk}}}-\sqrt{k-1})/(\tau\sqrt d)}, where \AtlasInlineMath{c_0,c_1>0} is the source universal constant.

\textbf{Quantity.} Q1

\textbf{Resource.} \AtlasInlineMath{N_{\mathrm{vec}}=b_{\mathrm{blk}}N_{\mathrm{block}}}

\textbf{Status.} Lower bound missing

\textbf{Lower bound.} No admitted bound for this comparison

\textbf{Upper bound.} B1

\AtlasMath{N_{\mathrm{block}}\le C\frac{\Lambda}{g}\log\frac{d\tau}{\varepsilon},\qquad N_{\mathrm{vec}}=b_{\mathrm{blk}}N_{\mathrm{block}}}
\textbf{Comparison argument.} An attainable call upper bound for a fixed algorithm. Minimum algorithmic complexity is not defined, so a missing lower bound is not a general minimax gap.

\needspace{7\baselineskip}\section{4. Bounds and supporting evidence}
The following conditions apply to every bound in this profile, in addition to each bound's stated region and algorithm scope.

\textbf{Oracle mapping.} The total-perturbation and U-projected-perturbation bounds hold jointly at every step. Neither independence nor zero means are used.

\textbf{Resource mapping.} L block calls of width b\_blk cost b\_blk L column queries.

\needspace{4\baselineskip}\subsection{B1 - Fixed Noisy Power Method with randomized initialization and causal bounded perturbations.}
\textbf{Admission.} Included in the bracket

\textbf{Direction.} upper

\textbf{Witness or lower-bound scope.} Randomized witness

\AtlasMath{N_{\mathrm{block}}\le C\frac{\Lambda}{g}\log\frac{d\tau}{\varepsilon},\qquad N_{\mathrm{vec}}=b_{\mathrm{blk}}N_{\mathrm{block}}}
\textbf{Sources.} \href{https://www.cs.utexas.edu/~ecprice/papers/npm.pdf}{[1]}

\textbf{Simultaneous resource costs.} All coordinates refer to this one witness or necessary constraint.

\textbf{N block.}

\AtlasMath{O((\Lambda/g)\log(d\tau/\varepsilon))}
\textbf{N vec.} b\_blk*N\_block

\textbf{Constants.} The iteration constant C and initialization-tail constants c\_0 and c\_1 are universal constants from the source theorem.

\textbf{Source location.} \href{https://www.cs.utexas.edu/~ecprice/papers/npm.pdf}{[1]} pdf pages: 1; 2; 3; theorem or section: Noisy Power Method; Corollary 1.1; theorem: Noisy Power Method; Corollary 1.1

\needspace{7\baselineskip}\section{References}
[1] Moritz Hardt; Eric Price. The Noisy Power Method: A Meta Algorithm with Applications. 2014. author manuscript dated November11,2014.

\url{https://www.cs.utexas.edu/~ecprice/papers/npm.pdf}

\end{document}
